\section*{Bab \ref{ch:b2:elements-in-the-earth} --- Unsur-Unsur di dalam Bumi}

\begin{solution}{exo:b2:elements-in-the-earth:1}
Kalium adalah $2\times39.10/(2\times39.10+16.00) = 78.20/94.20 = 0.8301$ bagian massa
\ce{K2O}, sehingga $\num{5130}\times0.8301 = \qty{4259}{ppm}$ K. Titanium adalah
$47.87/79.87 = 0.5993$ bagian \ce{TiO2}: $\num{27310}\times0.5993 = \qty{16370}{ppm}$ Ti.
\end{solution}

\begin{solution}{exo:b2:elements-in-the-earth:2}
Aluminium bersifat litofil dan refraktori: Bumi silikat mempertahankan seluruhnya, dalam proporsi yang hampir
kondritik (1.16). Natrium litofil tetapi volatil (\qty{958}{K}): unsur ini
hilang sebelum atau selama akresi, sehingga tersisa 0.22. Nikel mengembun bersama magnesium
(\qty{1353}{K}), sehingga volatilitas tidak dapat menjelaskan 0.077: nikel siderofil, dan sebagian besarnya
masuk ke inti. Iridium refraktori (\qty{1603}{K}) dan sangat siderofil:
0.0029, hampir seluruhnya berada di inti.
\end{solution}

\begin{solution}{exo:b2:elements-in-the-earth:3}
$\mathrm{La}_N = 31/0.232 = 133.6$, $\mathrm{Yb}_N = 2.0/0.163 = 12.27$, dan
$(\mathrm{La}/\mathrm{Yb})_N = 10.9$. Kerak diperkaya dalam semua tanah jarang, dan
sebelas kali lebih banyak dalam yang ringan daripada yang berat: kerak tersusun atas lelehan, yang
paling mudah mengambil tanah jarang ringan yang inkompatibel.
\end{solution}

\begin{solution}{exo:b2:elements-in-the-earth:4}
$\mathrm{Sm}_N = 4.7/0.145 = 32.4$, $\mathrm{Eu}_N = 1.0/0.0546 = 18.3$,
$\mathrm{Gd}_N = 4.0/0.198 = 20.2$; $\mathrm{Eu}^* = \sqrt{32.4\times20.2} = 25.6$ dan
$\mathrm{Eu}/\mathrm{Eu}^* = 18.3/25.6 = 0.72$, anomali negatif. Plagioklas, yang
menerima \ce{Eu^2+}, mengkristal dari lelehan yang membangun kerak atas dan tertinggal
lebih dalam.
\end{solution}

\begin{solution}{exo:b2:elements-in-the-earth:5}
$\bar D_{\mathrm{La}} = 0.65\times0.058+0.18\times0.001+0.17\times0.001 = 0.0377 +
0.0002 + 0.0002 = 0.038$; $\bar D_{\mathrm{Lu}} = 0.3445+1.044+0.0306 = 1.42$. Lutesium
kompatibel ($\bar D > 1$), lantanum sangat inkompatibel. Garnet memuat
$1.044/1.42 = 74$ persen lutesium batuan itu, meskipun garnet hanya 18 persen
batuan.
\end{solution}

\begin{solution}{exo:b2:elements-in-the-earth:6}
Untuk natrium, Na/Mg Bumi silikat adalah $2670/\num{228000} = 0.01171$ dan nilai CI
$5010/\num{95870} = 0.05226$: perbandingannya 0.224. Untuk nikel, $1960/\num{228000}
= 0.008596$ terhadap $\num{10640}/\num{95870} = 0.1110$: 0.0775. Natrium
terkuras karena volatilitasnya (\qty{958}{K}); nikel mengembun di atas \qty{1300}{K} dan
tiga kali lebih terkuras daripada natrium karena bersifat siderofil: inti mengambil sebagian besar
nikel.
\end{solution}

\begin{solution}{exo:b2:elements-in-the-earth:7}
La: $5.21\times0.058 = \qty{0.302}{ppm}$; Lu: $0.53\times0.53 = \qty{0.281}{ppm}$. Nilai
penormal saling meniadakan dalam perbandingan dari perbandingan:
$(\mathrm{La}/\mathrm{Lu})_{N,\mathrm{cpx}}/(\mathrm{La}/\mathrm{Lu})_{N,\mathrm{melt}} =
D_{\mathrm{La}}/D_{\mathrm{Lu}} = 0.058/0.53 = 0.11$: pola kristal itu sembilan
kali lebih curam ke arah tanah jarang berat daripada pola lelehan.
\end{solution}

\begin{solution}{exo:b2:elements-in-the-earth:8}
Kerak: $(12/0.658)/(31/0.648) = 18.24/47.84 = 0.38$. Basal Hawaii:
$(18.10/0.658)/(15.20/0.648) = 27.51/23.46 = 1.17$. MORB: $7.96/8.04 = 0.99$. Hanya
kerak, dengan kekurangan niobium yang kuat, yang membawa ciri zona subduksi.
\end{solution}

\begin{solution}{exo:b2:elements-in-the-earth:9}
\omterm{def:b2:elements-in-the-earth:bse}{Mantel primitif}: $240/0.0203 = \num{11800}$. MORB: K $= 0.160\times0.830\times10^4 =
\qty{1328}{ppm}$, K/U $= \num{11200}$. Kerak: K $= 2.80\times0.830\times10^4 =
\qty{23240}{ppm}$, K/U $= \num{8600}$. Kalium naik seratus kali lipat dari mantel ke kerak,
tetapi K/U berubah kurang dari sepertiga: kedua unsur kira-kira sama inkompatibelnya,
dan peleburan hampir tidak memisahkan keduanya. Uranium refraktori, sehingga kelimpahannya mengikuti
model kondritik; K/U yang diukur pada lava kemudian memberikan kalium planet itu, yang,
karena volatil, tidak dapat diambil dari \omterm{def:b2:planet-formation-meteorites:chondrule}{kondrit}.
\end{solution}

\begin{solution}{exo:b2:elements-in-the-earth:10}
Dalam kesetimbangan $C_{\mathrm{rock}} = \bar D\,C_{\mathrm{melt}}$ untuk setiap unsur (dan
fraksi lelehan yang kecil hampir tidak mengubah komposisi batuan), sehingga
$(C_1/C_2)_{\mathrm{melt}} = (C_{1}/\bar D_1)/(C_{2}/\bar D_2)$ untuk $C_1,
C_2$ batuan, dan perbandingannya adalah $\bar D_2/\bar D_1$. Untuk La dan Lu: $1.419/0.03805 = 37$. Lelehan
mantel yang mengandung garnet 37 kali lebih kaya La relatif terhadap Lu daripada sumbernya: garnet
menahan tanah jarang berat, dan pola lelehan itu curam.
\end{solution}

\begin{solution}{exo:b2:elements-in-the-earth:11}
$C_{\mathrm{mix}} = fC_U + (1-f)C_M$ untuk setiap unsur: linear terhadap $f$. Perbandingan
$\bigl(f\,\mathrm{Th}_U + (1-f)\mathrm{Th}_M\bigr)/\bigl(f\,\mathrm{Nb}_U +
(1-f)\mathrm{Nb}_M\bigr)$ adalah hasil bagi dua fungsi linear, sebuah hiperbola dalam $f$.
Th/Nb $= 0.2$ ketika $f(10.5-2.4) + (1-f)(0.404-1.048) = 0$, yaitu $8.1f = 0.644(1-f)$,
$f = 0.644/8.744 = 0.074$: sekitar 7 persen kerak menaikkan Th/Nb dari 0.077 menjadi 0.2. Untuk
$f = 0.1$: Th $= 1.05 + 0.364 = 1.414$, Nb $= 1.2 + 4.716 = 5.916$, Th/Nb $= 0.24$.
\end{solution}

\begin{solution}{exo:b2:elements-in-the-earth:12}
Membuang komponen yang tidak memuat kedua unsur itu tidak mengubah massa keduanya, dan karena itu perbandingannya.
Jika seluruh aluminium dipertahankan, $(\mathrm{Al}/\mathrm{Mg}) = 1.16$ kali CI
berarti Bumi silikat memuat $1/1.16 = 0.86$ magnesium kondritiknya: 14 persen
hilang. Silikon relatif terhadap aluminium bernilai $0.83/1.16 = 0.72$ kali nilai CI:
28 persen silikon kondritik akan hilang.
\end{solution}

\begin{solution}{pb:b2:elements-in-the-earth:1}
\textbf{1.} Fe/Mg $= \num{182800}/\num{95870} = 1.907$; Ni/Mg $= 0.1110$; Si/Mg $=
1.111$; Al/Mg $= 0.0887$.

\textbf{2.} $(\num{182800}/55.85)/(\num{106500}/28.09) = 3273/3791 = 0.863$ atom besi per
atom silikon, dalam selisih 7 persen dari 0.81 fotosfer: masuk akal, karena \omterm{def:b2:elements-in-the-earth:ci}{kondrit CI}
sesuai dengan Matahari untuk unsur-unsur pembentuk batuan.

\textbf{3.} Hilangnya komponen yang tidak memuat kedua unsur suatu perbandingan tidak mengubah
perbandingan itu (\cref{exo:b2:elements-in-the-earth:12}): hanya konsentrasi mutlaknya yang berubah.

\textbf{4.} Magnesium litofil dan tidak masuk ke inti; magnesium adalah unsur utama,
yang kadarnya diketahui baik dalam \omterm{def:b2:planet-formation-meteorites:chondrule}{kondrit} maupun dalam mantel. Besi adalah unsur yang terbagi
antara inti dan mantel: besi tidak dapat menjadi acuannya sendiri.

\textbf{5.} Unsur-unsur yang mengembun bersama tidak terpisahkan oleh volatilitas di
nebula, sehingga Bumi total semestinya mempertahankan perbandingannya. Al/Mg menyimpang: aluminium
refraktori (\qty{1653}{K}), dan Al/Mg Bumi silikat bernilai 1.16 kali nilai CI.

\textbf{6.} Fe $= 8.05\times55.85/71.85 = 6.26$ persen; Ni $= 0.25\times58.69/74.69 =
0.196$ persen.

\textbf{7.} Fe/Mg $= 6.26/22.8 = 0.2746$; Ni/Mg $= 0.196/22.8 = 0.00860$.

\textbf{8.} Besi: $1 - 0.2746/1.907 = 0.856$, 85.6 persen hilang; nikel:
$1 - 0.00860/0.1110 = 0.923$, 92.3 persen.

\textbf{9.} Iridium lebih kuat siderofil daripada nikel, dan nikel daripada besi:
makin kuat suatu unsur menyukai logam, makin banyak bagiannya yang dibawa logam ke inti.

\textbf{10.} Dalam silikat, sebagai \ce{Fe^2+} dalam olivin dan piroksen (8 persen FeO):
besi yang teroksidasi menjadi FeO berikatan dengan oksigen seperti \omterm{def:b2:elements-in-the-earth:goldschmidt}{unsur litofil} dan tetap berada di
mantel; hanya besi logam yang tenggelam.

\textbf{11.} $(1-x)C_{\mathrm{Fe}} + xc_{\mathrm{Fe}} = r_{\mathrm{Fe}}(1-x)C_{\mathrm{Mg}}$;
membagi dengan $(1-x)c_{\mathrm{Fe}}$ menghasilkan proposisi itu.

\textbf{12.} $r_{\mathrm{Fe}}C_{\mathrm{Mg}} - C_{\mathrm{Fe}} = 1.907\times22.8 - 6.26 =
37.22$; $y = 37.22/85.5 = 0.4353$, $x = y/(1+y) = 0.303$.

\textbf{13.} $0.1110\times22.8 - 0.196 = 2.335$; $y = 2.335/5.2 = 0.4490$, $x = 0.310$.

\textbf{14.} $y = (85.5\times37.22 + 5.2\times2.335)/(85.5^2 + 5.2^2) =
3194.4/7337.3 = 0.4354$, $x = 0.303$.

\textbf{15.} Setiap persamaan berbobot sebesar kuadrat koefisiennya: $85.5^2/5.2^2 = 270$
kali lebih besar untuk besi, unsur yang jauh lebih melimpah.

\textbf{16.} $37.22/2.335 = 15.9$, terhadap $85.5/5.2 = 16.4$ pada model: selisih 3
persen.

\textbf{17.} $0.303\times\num{5.972e24} = \qty{1.81e24}{kg}$.

\textbf{18.} $(32.5-30.3)/32.5 = 7$ persen lebih rendah.

\textbf{19.} $y = 0.325/0.675 = 0.4815$, Fe dalam inti $= 37.22/0.4815 = 77.3$ persen,
dan Ni $= 2.335/0.4815 = 4.85$ persen: 18 persen unsur ringan, kira-kira dua kali 5 hingga
10 persen yang diizinkan massa jenis. Fe/Mg Bumi pasti agak lebih tinggi daripada nilai
CI, atau magnesium silikatnya lebih tinggi daripada nilai model.

\textbf{20.} Silikon, belerang, karbon, oksigen, dan hidrogen. Silikon kurang dalam
Bumi silikat: perbandingan ternormalisasi-Mg-nya, 0.83, berada di bawah tingkat refraktori.

\textbf{21.} $1.111\times22.8 - 21.0 = 4.331$; $4.331/0.4354 = 9.9$ persen Si dalam
inti.

\textbf{22.} 9.9 terhadap 6 persen: inti mungkin memuat lebih banyak silikon daripada model, atau
sebagian silikon tidak pernah terakresi. Silikon mengembun pada \qty{1310}{K}, sedikit
di bawah magnesium dan besi (1336 dan \qty{1334}{K}): sebagian mungkin tetap berada dalam gas.

\textbf{23.} Inti: $0.325\times0.5 = \qty{0.1625}{ppm}$ massa Bumi; Bumi
silikat: $0.675\times0.001 = \qty{0.000675}{ppm}$. Inti memuat
$0.1625/0.1632 = 99.6$ persen emas.

\textbf{24.} Kimia saja menempatkan
\textbf{\qty[text-series-to-math=true]{30.3}{\percent}} massa Bumi di intinya,
dalam selisih 7 persen dari 32.5 persen hasil seismologi.
\end{solution}
